\documentclass[a4paper,12pt]{article}

\usepackage{graphicx}
\usepackage{hyperref}
\usepackage{listings}
\usepackage{xcolor}

\title{Arduino LED Blink Project}
\author{Alireza Naderinasab}
\date{\today}

\begin{document}
	
	\maketitle
	
	\section{Project Description}
	This project blinks an LED connected to pin 11 every 1 second using Arduino IDE.
	
	\section{Hardware Required}
	\begin{itemize}
		\item Arduino Uno
		\item LED
		\item 100$\Omega$ resistor
		\item Breadboard and wires
	\end{itemize}
	
	\section{Circuit Connection}
	LED is connected to digital pin 11 of Arduino.
	
	\section{Code}
	
	\begin{lstlisting}[language=C]
		int ledPin = 11;
		
		void setup() {
			pinMode(ledPin, OUTPUT);
		}
		
		void loop() {
			digitalWrite(ledPin, HIGH);
			delay(1000);
			
			digitalWrite(ledPin, LOW);
			delay(1000);
		}
	\end{lstlisting}
	
	\section{How to Run}
	\begin{enumerate}
		\item Open Arduino IDE
		\item Upload blink.ino to Arduino board
		\item LED will start blinking
	\end{enumerate}
\section{Simulation Output}

\subsection{LED ON State}

\begin{center}
	\includegraphics[width=0.6\textwidth]{images/On.jpg}
\end{center}

\subsection{LED OFF State}

\begin{center}
	\includegraphics[width=0.6\textwidth]{images/Off.jpg}
\end{center}
\section{Conclusion}

In this project, a simple LED blinking system was implemented using Arduino IDE. 
The system uses digital pin control to switch the LED ON and OFF with a fixed delay.

The project demonstrates the basic principles of embedded systems, including GPIO control, timing functions, and hardware interaction.

This implementation can be extended in future to more complex systems such as sensor-based automation or IoT applications.
\end{document}